\documentclass[11pt]{article}
\usepackage[margin=1in]{geometry}
\usepackage{amsmath,amssymb,booktabs,longtable,array,enumitem,hyperref}
\newcommand{\E}{\mathcal E}
\newcommand{\Proc}{\mathsf{Proc}_{\mathrm{LF}}}
\title{V21 P4-T03: Emergence-First Locally Finite Process-Order Option}
\author{Ontology Formalizer parent--child synthesis}
\date{2026-07-21}

\begin{document}
\maketitle

\begin{center}
\fbox{\parbox{0.93\linewidth}{\textbf{Status and authority.}
This is a \textbf{draft/control, proposal-only} ontology option and
\textbf{source-extension data}. It is not a canonical-ontology candidate,
adoption, rejection, Gate Chair decision, continuum-emergence result, proof of
GR, or publication artifact. Current status is
\texttt{blocked\_adoption\_open\_continuation}.}}
\end{center}

\section{Question and exact boundary}

The bounded question is whether one can state a coherent source option that is
genuinely weaker than a differentiable four-dimensional manifold, so that
dimension and continuum structure would have to be reconstructed. The answer
below supplies source syntax, finite controls, and reconstruction obligations.
It does not show that any continuum emerges or that the source is physical.

No topology, dimension, coordinate, differentiability, atlas, metric, volume,
physical time, causal cone, clock, matter field, coupling, action, probability,
or gravitational equation is admitted as source data. The order relation is
formal source precedence only. P4-T03 does not compare, select, adopt, or reject
an ontology regime; those remain later and protected questions.

\section{Minimum source signature}

A proposal instance is
\[
  \mathfrak O_{\mathrm{EF}}^{\mathrm{LF}}=(E,\prec),
\]
where $E$ is a nonempty finite or countably infinite set of source tokens and
$\prec\,\subseteq E\times E$ satisfies:
\begin{enumerate}[label=(S\arabic*)]
  \item irreflexivity: $x\not\prec x$ for every $x\in E$;
  \item transitivity: $x\prec y\prec z$ implies $x\prec z$;
  \item local finiteness: if $x\prec y$, then
  $I(x,y)=\{z\in E:x\prec z\prec y\}$ is finite.
\end{enumerate}
The reflexive closure, finite intervals, antichains, minimal and maximal
elements, and finite-order isomorphisms are derived syntax. Labels,
probabilities, amplitudes, action values, coordinates, tensor types, units,
clocks, detectors, and matter species are not in the minimum signature.
The word ``event'' means only a neutral source token.

\section{Finite process composition}

A process fragment $P=(F,\prec_F,i,o)$ consists of a finite strict partial
order $(F,\prec_F)$ and injections from finite interface sets
$i:B_{\mathrm{in}}\to F$ and $o:B_{\mathrm{out}}\to F$. The input image is
an antichain of minimal elements and the output image is an antichain of
maximal elements. Interfaces are algebraic boundaries, not spacelike
hypersurfaces or simultaneity classes.

Given an interface bijection
$\phi:B_{\mathrm{out}}(P_1)\to B_{\mathrm{in}}(P_2)$, form the disjoint union,
identify $o_1(b)$ with $i_2(\phi(b))$, and take the transitive closure of the
two internal orders. The composite is defined only if the quotient relation is
irreflexive. Identity fragments have carrier equal to the interface and empty
strict order. On the domain where all intermediate composites are defined,
iterated quotient and closure is associative up to source-order isomorphism.
This is a finite algebraic composition certificate, not a physical dynamics,
continuum limit, or clock evolution.

A countable history may be represented as a compatible directed union of
finite fragments, but local finiteness of that union must be checked; it does
not follow automatically from finite-stage admissibility.

\section{Strict-weakness certificate}

\begin{longtable}{p{0.23\linewidth}p{0.29\linewidth}p{0.39\linewidth}}
\toprule
Layer & Present in source & Absent and therefore still a burden \\
\midrule
\endhead
Carrier & token equality and membership & topology, neighborhoods,
connectedness, compactness, spatial slices \\
Precedence & one locally finite strict order & physical time, causal cones,
orientation, duration, Lorentzian interval \\
Composition & finite order interfaces and guarded gluing & chart overlaps,
smooth transitions, field equations, transition probabilities \\
Geometry & none & dimension, coordinates, differentiability, tangent spaces,
metric, signature, connection, curvature, volume, scale \\
Physics & none & clocks, observers, detectors, matter, coupling, stress energy,
action, Einstein dynamics, empirical semantics \\
\bottomrule
\end{longtable}

Thus the declared primitives are strictly weaker than target manifold
structure by direct signature inventory. This is a syntactic certificate, not
a theorem that the missing structures are reconstructible or unique.

\section{Finite controls and target-dimension nonselection}

The four-token diamond has covers
$a\prec b$, $a\prec c$, $b\prec d$, and $c\prec d$, with $b$ and $c$
incomparable. It composes two fragments along the interface $\{b,c\}$.
A second control has two disjoint two-token chains $p\prec q$ and
$r\prec s$; it shows that connectedness and spatial interpretation are not
source axioms.

For an evaluator-side illustration only, embed the diamond into Minkowski
targets $T_2=(\mathbb R^2,\eta_2)$ and $T_4=(\mathbb R^4,\eta_4)$ by
\[
 a\mapsto(-2,0,\ldots,0),\quad
 b\mapsto(0,-\tfrac12,0,\ldots,0),\quad
 c\mapsto(0,\tfrac12,0,\ldots,0),\quad
 d\mapsto(2,0,\ldots,0).
\]
The strict chronological relation induced on the four images is the same
diamond in both target dimensions. This finite witness shows only that
order-fidelity for one finite fragment does not select dimension. It does not
satisfy a continuum limit, counting law, stability, uniqueness, dynamics, or
physical-causality burden.

\section{Unfilled dynamics target}

A future source law could assign a transition kernel $K(P\to P')$ or amplitude
$A(P\to P')$ to admissible fragment extensions. Such a law would have to be
invariant under source-order isomorphism, state its normalization or
composition rule, satisfy finite-stage consistency, and be defined without a
target atlas, target metric, target dimension, physical clock, or GR dynamics.
No kernel, amplitude, action, history measure, growth rule, or selection rule
is supplied or adopted here.

\section{Target-only realization obligation}

Let $T=(M,g,o_t,\triangleleft_T,\mathcal R,\rho,L,\epsilon)$ be an entirely
target-side test object, and let $\iota:E\to M$ be an evaluator-side injection.
A future realization predicate may test order fidelity, local finiteness in
declared target regions, coarse number--volume agreement, intrinsic recovery,
stability, and uniqueness. Every entry of $T$ and the map $\iota$ remain
outside $\mathfrak O_{\mathrm{EF}}^{\mathrm{LF}}$. They may not define source
admissibility, source dynamics, or a manifoldlikeness selector.

A lawful continuum claim would require a source-internal family and limit,
target-neutral manifoldlikeness discriminator, topology recovery, convergent
dimension estimates, recovered differentiability, a physical-causality bridge,
measure and scale calibration, metric recovery, uniqueness and stability,
dynamical typicality, operational semantics, matter coupling, and GR recovery.
The machine-readable burden map records these eleven layers. No current finite
example discharges any of them.

\section{Relation to current project ontology and exact GR}

The registered canonical ontology currently uses a four-dimensional
differentiable substrate as interpretive framing. This task-local option does
not edit, reject, supersede, or silently reinterpret that authority. It is an
alternative proposal for later comparison only. The exact-GR benchmark remains
target-side; its metric, causal structure, matter content, equations, and
empirical success cannot serve as source premises.

Current ontology does not derive a manifoldlikeness selector, continuum limit,
source dynamics, or reconstruction chain from the proposed process-order
signature. That statement is not a no-go theorem: a conservative source-side
extension remains open, so adoption is blocked while continuation remains open.

\section{Disposition}

P4-T03 is satisfied at proposal scope by one coherent minimum source option,
guarded composition calculus, two finite controls, a target-dimension
nonselection witness, and an explicit manifoldlikeness burden map. The
scientific Distance-to-GR ledger is unchanged because no source package,
dynamics, reconstruction law, continuum, metric, matter model, coupling, or
ontology regime is adopted. P4-T04 is now dependency-ready for a separate
comparative AgentJob. P4-T05 remains human-gated.

\section*{Primary literature context}

Locally finite partial order is used here only as a modeling precedent from
Bombelli, Lee, Meyer, and Sorkin (1987). Sequential growth is cited only as an
example that dynamics is a separate additional law (Rideout \& Sorkin, 2000).
Continuum-topology recovery is cited only as evidence that reconstruction
requires nontrivial assumptions and tests (Major, Rideout, \& Surya, 2007).
These sources do not establish this project-specific proposal, physical
causality, continuum emergence, ontology adoption, or GR recovery.

\begin{enumerate}
  \item Bombelli, L., Lee, J., Meyer, D., \& Sorkin, R. D. (1987).
  Space-time as a causal set. \emph{Physical Review Letters, 59}(5), 521--524.
  \url{https://doi.org/10.1103/PhysRevLett.59.521}
  \item Rideout, D. P., \& Sorkin, R. D. (2000). A classical sequential
  growth dynamics for causal sets. \emph{Physical Review D, 61}, 024002.
  \url{https://doi.org/10.1103/PhysRevD.61.024002}
  \item Major, S., Rideout, D., \& Surya, S. (2007). On recovering continuum
  topology from a causal set. \emph{Journal of Mathematical Physics, 48},
  032501. \url{https://doi.org/10.1063/1.2435599}
\end{enumerate}

\end{document}
